\subsection*{Примитивные объекты и их представление в памяти}
\begin{enumerate}
    \item \textbf{Точка:} пара значений $(x, y)$.
    \item \textbf{Вектор:} пара значений $(x, y)$. По сути, отображают представление вектора в ортонормированном базисе.\newline
    Вектор $\Vec{AB}$ можно задать как пару значений $(x_b - x_a, y_b - y_a)$
    \item \textbf{Прямая:} Можно представить 3 основными способами: \begin{enumerate}
        \item Тройка значений $(A, B, C)$, задающих уравнение прямой $Ax+By+C = 0$.\newline
        В этом случае, нормаль-вектор равен $(A, B)$, направляющий вектор прямой равен $(-B, A)$
        \item Две точки, которые образуют прямую. В таком случае, согласно каноническому уравнению прямой \begin{equation*}
        \frac{x - x_2}{x_1 - x_2} = \frac{y - y_2}{y_1 - y_2}
        \end{equation*}
        можно выразить тройку $(A, B, C)$
        \item Точка на прямой и направляющий вектор.
    \end{enumerate}
    \item \textbf{Луч:} точка-начало луча и направляющий вектор. Его, как и прямую, можно представить параметрически: $\Vec{r} = \Vec{r_0} + \Vec{a}t$, но $t > 0$.
    \item \textbf{Отрезок:} пара точек - концы отрезка.
    \item \textbf{Окружность:} точка-центр окружности $(x_0, y_0)$ и радиус $R$. \newline
    Уравнение окружности: $(x - x_0)^2 + (y - y_0)^2 = R^2$
\end{enumerate}

\subsection*{Основные операции}
\begin{enumerate}
    \item \Def Пусть $v = (x_1, y_1)$ и $u = (x_2, y_2)$ — элементы $\R^2$. Тогда \textbf{cкалярное произведение} $(a, b) = x_ax_b + y_ay_b$. Известно, что $(a, b) = |a|\cdot|b|\cdot \cos(a \textasciicircum b)$
    \item \Def Пусть $a = (x_a, y_a)$ и $b = (x_b, y_b)$ — элементы $\R^2$. Тогда \textbf{косое} или \textbf{псевдовекторное произведение} $[a, b] = x_ay_b - x_by_a = \begin{vmatrix}
        x_a & y_a \\
        x_b & y_b 
    \end{vmatrix}$. Известно, что $[a, b] = |a|\cdot|b|\cdot sin (a \textasciicircum b)$
    \item \textbf{Длина вектора}\textit{(он же длина отрезка)}. $|\Vec{a}| = \sqrt{(x_2 - x_1)^2 + (y_2 - y_1)^2}$
\end{enumerate}

\subsection*{Взаимное расположение объектов}
\begin{enumerate}
    \item \textbf{Точка + точка}, он же отрезок. Длина отрезка $AB = \sqrt{(x_2 - x_1)^2 + (y_2 - y_1)^2}$
    \item \textbf{Точка + прямая}

    Пусть дана точка $A(x_1, y_1)$ и прямая $l: Ax + By + C = 0$. Тогда, точка лежит на прямой, если $Ax_1 + By_1 + C = 0$.
    
    \textbf{Расстояние от точки до прямой} равно: $\displaystyle \rho(A, l) = \frac{|l(A)|}{\sqrt{A^2 + B^2}}$

    \drawsomesmall{pix/dot_line.jpg}
    
    \item \textbf{Точка + луч} 

    Пусть даны точка $P$ и луч $k$ с началом в точке $A$ и направляющим вектором $\Vec{a}$. 
    
    Тогда, $P \in k \Leftrightarrow \begin{cases}
    (\Vec{a}, \Vec{AP}) > 0 \textit{(значит вектора сонаправлены)} \\
    [\Vec{a}, \Vec{AP}] = 0\textit{(значит вектора коллинеарны)}.
    \end{cases}
    $

    \textbf{Расстояние от точки до луча.}
    
    $\rho(A, k) = \begin{cases}
    \rho(A, l(k)), & (\Vec{a}, \Vec{AP}) > 0 \\
    |PA|, & \textit{иначе}.
    \end{cases}
    $, 
    
    где $l(k)$ - прямая, содержащая луч $k$.

    \drawsomesmall{pix/dot_ray.jpg}
    \item \textbf{Точка + отрезок}

    Если точка $C$ лежит на отрезке $AB$, то площадь треугольника, который они образуют, должна быть равна 0. Кроме того, векторы от точки до концов отрезка должны быть противоположно направлены, иначе эта точка лежит вне отрезка. Итого, $C \in AB \Leftrightarrow \begin{cases}
        [\Vec{CA}, \Vec{CB}] = 0 \\
        (\Vec{CA}, \Vec{CB}) < 0
    \end{cases}$ 

    \textbf{Расстояние от точки до отрезка.} Если проекция точки $C$ лежит на отрезке(т. е. углы в треугольнике острые: $\Rightarrow (\Vec{AC}, \Vec{AB}) \cdot (\Vec{BC}, \Vec{BA}) > 0$), тогда расстояние можно найти как расстояние до прямой, содержащей отрезок, а иначе - нужно взять минимум из расстояний до концов отрезка. Итого, $\rho(C, AB) = \begin{cases}
        \rho(C, l(AB)), & pr_{AB}C \in AB \\
        min(|CA|, |CB|), & \textit{иначе}
    \end{cases}$

    \drawsomesmall{pix/dot_segment.jpg}
    \item \textbf{Точка + окружность} 
    
    Точка $A(x_1, y_1)$ лежит на окружности $(O, R)$, если она удовлетворяет уравнению окружности. Расстояние от точки до окружности - $|R - OA|$.
    \item \textbf{Прямая + прямая} 

    Две прямые $a_1:A_1x+B_1y+C_1=0$ и $a_2:A_2x+B_2y+C_2=0$ параллельны, если их направляющие вектора или нормаль-векторы коллинеарны. Прямые совпадают, если у них помимо этого, совпадает еще и свободный член в уравнении. Расстояние между двумя параллельными прямыми равно расстоянию от любой точки с одной прямой до другой прямой.

    \textbf{Точка пересечения двух непараллельных прямых.} Её можно найти, решив систему уравнений $\begin{cases}
        A_1x+B_1y+C_1=0\\
        A_2x+B_2y+C_2=0
    \end{cases}$
    
    \item \textbf{Прямая + луч} Очевидно, что если направляющие векторы луча и прямой коллинеарны, то они не пересекаются. Если же при этом точка-начало луча лежит на прямой, то луч лежит на этой прямой.

    Теперь рассмотрим случай, когда направляющие векторы не коллинеарны. Для этого нужно решить систему уравнений $\begin{cases}
        Ax+By+C=0\\
        x = x_0 + a_xt\\
        y = y_0 + a_yt
    \end{cases} \Rightarrow t = \displaystyle\frac{-Ax_0-By_0-c}{Aa_x + Ba_y}$. Если $t \Ge 0$, тогда луч и прямая действительно пересекаются. А если нет, то лучи и прямая не имеют общих точек.

    \textbf{Точка пересечения прямой и луча.} Если $t \Ge 0$, то точку пересечения можно найти просто подставив это значение в уравнение луча. Если же $t < 0$, то существует только \textit{мнимая} точка пересечения.
    
    \item \textbf{Прямая + отрезок}

    Прямая $l$ пересекает отрезок $PQ$ только в том случае, когда концы отрезка находятся по разные стороны от прямой. Т.е. $l(P)\cdot l(Q) \Le 0$.

    \textbf{Расстояние от прямой до отрезка.} Очевидно, что если прямая и отрезок пересекаются, то расстояние между ними равно 0, а если параллельны, то  расстояние равно расстоянию от любой точки отрезка до прямой.

    Если же нет, тогда можно заметить, что $\rho(PQ, l) = \min(\rho(P, l), \rho(Q, l))$. Это можно доказать, если, например, продлить отрезок в прямую, и посмотреть как меняется расстояние между прямыми при движении по одной из них.
    \drawsomemedium{pix/line_ray.jpg}
    
    \item \textbf{Прямая + окружность}

    Прямая $l$ пересекает окружность $(O, R)$ в том случае, когда $\rho(O, l) \Le R$.

    \textbf{Точки пересечения прямой и окружности.} Пусть $O^{'} = pr_{l}O$, тогда, расстояние от $O^{'}$ до точек пересечения $A, B$ равно $d = \sqrt{R^2 - {OO^{'}}^2}$.
    
    Тогда, точки $A, B$ можно получить как $\displaystyle O^{'} \pm d \cdot \frac{\Vec{a}}{|a|}$.

    \drawsomesmall{pix/line_circle.jpg}
    
    \item \textbf{Луч + Луч}

    Пусть у нас есть луч 1 и луч 2. Возьмем вместо луча 1 прямую, на которой он лежит, и проверим, пересекает ли ее луч 2. Это мы умеем делать. Если нет, то раз луч 2 не пересекает прямую, в которой лежит луч 1, то и сам луч 1 он, очевидно, тоже не пересекает.
    
    Теперь возьмем наоборот. Если луч 1 пересекает прямую, на которой лежит луч 1, то значит и оба луча пересекаются. Заметим, что без этой проверки нельзя: мог быть случай, когда луч 2 пересекал луч 1 на \textit{мнимой} части.

    \item \textbf{Луч + отрезок}

    Если луч пересекает отрезок, тогда: $\begin{cases}
        l(P) \cdot l(Q) \Le 0 \\
        \angle PAD + \angle DAQ \Le \pi,
    \end{cases}$

    где $l$ - прямая, на которой лежит луч. 
    
    Последнее утверждение нужно для того, чтобы точки отрезка находились после точки-начала

    На центральном и правом изображениях приведены примеры, когда нарушается каждый из пунктов соответственно.
    \drawsomemedium{pix/ray_segment.jpg}
    \item \textbf{Луч + окружность} 
    Луч пересекает окружность, если ($(\Vec{OA}, \Vec{a}) \Ge 0$ и $\rho(O, \Vec{a}) \Le R$) или ($\rho(O, A) \Le R$)(если начало луча лежит в окружности, то $\forall$ луч будет пересекать окружность)

    \drawsomemedium{pix/ray_circle.jpg}

    \item \textbf{Отрезок + отрезок}

    Для того, чтобы проверить, пересекаются ли два отрезка нужно проверить, пересекает ли каждый из отрезков прямую, на которой лежит другой отрезок. То есть, $\begin{cases}
        AB \cap l(CD) \neq \varnothing \\
        CD \cap l(AB) \neq \varnothing
    \end{cases}$.

    \textbf{Расстояние между двумя отрезками.} Если они пересекаются, то расстояние равно 0. Если же нет, то $\rho(AB, CD) = \min(\rho(A, CD), \rho(B, CD), \rho(C, AB), \rho(D, AB))$.

    \drawsomemedium{pix/segment_segment.jpg}
    
    \item \textbf{Отрезок + окружность}
    
    Отрезок $AB$ пересекает окружность $(O, R)$, если $\begin{cases}
        \rho(O, AB) \Le R \\
        \begin{compound}
        |OA|^2 \Ge R^2 \\
        |OB|^2 \Ge R^2
        \end{compound}
    \end{cases}$
    
    \drawsomesmall{pix/segment_circle.jpg}
    
    \item \textbf{Окружность + окружность}

    Пусть есть две окружности $(O_1, R_1)$ и $(O_2, R_2)$. Для нахождения точек их пересечения решим систему $\begin{cases}
        (x - x_1)^2 + (y - y_1)^2 = R_1^2 \\
        (x - x_2)^2 + (y - y_2)^2 = R_2^2
    \end{cases}$

    Раскрыв скобки, получим выражение $2x(x_2-x_1) + 2y(y_2-y_1) + x_1^2 - x_2^2 + y_1^2 - y_2^2 = R_1^2 - R_2^2$, которое является уравнением прямой. После этого несложно найти точки их пересечения.
\end{enumerate}
\pagebreak